% TOP_PROBLEM: {"id": 20000921, "title": "Reverse Banaszczyk minimization and the ellipsoidal regime", "category_id": 2, "category": {"id": 2, "name": "combinatorics", "display_name": "Combinatorics", "description": "Counting problems, graph theory, discrete structures.", "slug": "combinatorics", "order_index": 2, "created_at": "2026-07-31T15:26:25.671Z"}, "status": "partially_solved", "problem_number": "AIM-COMBINATORICS-0046", "set_id": 14}
% FIRST_POSTED: 2026-10-01
% ENTRY_KIND: statement_only
% UnsolvedMath ID 20000921; source catalogue code AIM-COMBINATORICS-0046
% !TeX program = lualatex
% Definition and sources only; this file is not an LLM solution attempt.
\documentclass[11pt,a4paper]{article}
\usepackage[margin=25mm]{geometry}
\usepackage{fontspec}
\setmainfont{lmroman10-regular.otf}[BoldFont=lmroman10-bold.otf,
 ItalicFont=lmroman10-italic.otf,BoldItalicFont=lmroman10-bolditalic.otf]
\usepackage{amsmath,amssymb,mathtools}
\usepackage{unicode-math}
\setmathfont{latinmodern-math.otf}
\AtBeginDocument{\let\setminus\smallsetminus}
\usepackage{xurl,hyperref}
\hypersetup{colorlinks=true,urlcolor=blue}
\setcounter{secnumdepth}{0}
\setlength{\parindent}{0pt}
\setlength{\parskip}{5pt}
\setlength{\emergencystretch}{3em}
\begin{document}
\section*{Reverse Banaszczyk minimization and the ellipsoidal regime}\label{20000921}
{\small UnsolvedMath ID 20000921 · AIM-COMBINATORICS-0046 · Combinatorics · AIM Workshop Problem Lists}
\subsection{Definitions and mathematical statement}
Fix a dimension $d\geq1$ and a positive length bound $c$.
For a Borel set $K\subseteq\mathbb R^d$, its standard Gaussian volume is
\[
 \gamma_d(K)=(2\pi)^{-d/2}\int_K e^{-\|x\|_2^2/2}\,dx.
\]
A convex body here means a compact convex set with nonempty interior.
Say that $K$ has property $B_c$ if for every integer $m\geq0$ and every sequence $v_1,\ldots,v_m\in\mathbb R^d$ with $\|v_i\|_2\leq c$, there are signs $\epsilon_1,\ldots,\epsilon_m\in\{-1,1\}$ for which
\[
 \sum_{i=1}^m\epsilon_i v_i\in K.
\]
For $m=0$ the sum is $0$, so the property includes $0\in K$.
Sequences may repeat vectors, and the signing may depend on the entire sequence.

Determine the extremal quantity
\[
 b_d(c)=\inf\{\gamma_d(K):K\text{ is a convex body with }B_c\},
\]
and identify bodies attaining, or approaching, this infimum.
The source treats $c$ as a fixed small absolute constant; displaying it makes the dependence on the vector normalization explicit.
One may then ask how the smallest volume behaves with dimension for that fixed choice of $c$.
The universal requirement concerns every sequence length, not just a prescribed number of vectors.
No central symmetry is imposed unless one explicitly studies that additional restricted variant; a signed-sum requirement alone does not define symmetry of the body.
\subsection{Short English statement}
How small can the Gaussian volume of a convex body be if it can receive a signed sum of every bounded sequence of vectors?
\subsection{Sources}
\begin{itemize}
\item[S1] ULAM.AI and the UnsolvedMath Contributors, supplied problems.json, record 20000921 (AIM-COMBINATORICS-0046), AIM Workshop Problem Lists. Catalogue identity and source transcription; CC BY 4.0.
\url{https://huggingface.co/datasets/ulamai/UnsolvedMath}.
\item[S2] American Institute of Mathematics, Hereditary discrepancy and factorization norms, item 1.8. Original workshop question underlying AIM-COMBINATORICS-0046.
\url{http://aimpl.org/hereddiscrep/1/}.
\end{itemize}
\end{document}
